\DocumentMetadata{
  pdfversion=2.0,
  pdfstandard=ua-2,
  lang=en-US,
  tagging=on,
  tagging-setup={math/setup=mathml-SE,tabsorder=structure}
}
\documentclass[11pt,letterpaper]{article}
\usepackage[margin=0.68in,includefoot]{geometry}
\usepackage{newcomputermodern}
\usepackage{amsmath,amscd,mathtools}
\usepackage{tikz,adjustbox}
\providecommand{\square}{\mdlgwhtsquare}
\usepackage[hidelinks]{hyperref}
\hypersetup{
  pdftitle={Algebra First-Year Exam, January 2023, Part 2},
  pdfauthor={University of Florida Department of Mathematics},
  pdfsubject={Graduate mathematics examination},
  pdfdisplaydoctitle=true,
  bookmarksopen=true
}
\setcounter{secnumdepth}{0}
\setlength{\parindent}{0pt}
\setlength{\parskip}{0.28em}
\setlength{\emergencystretch}{3em}
\setlength{\leftmargini}{2.15em}
\setlength{\leftmarginii}{2.65em}
\setlength{\leftmarginiii}{3em}
\pagestyle{empty}
\raggedbottom
\allowdisplaybreaks
\NewTaggingSocketPlug{block/list/label}{examproblem}{%
  \tagstructbegin{tag=Lbl}\tagstructend%
  \tagstructbegin{tag=\LBody}%
  \tagstructbegin{tag=\ProblemHeadingTag,title-o={\ProblemHeadingText}}%
  \tagmcbegin{tag=\ProblemHeadingTag,actualtext={\ProblemHeadingText}}%
  #2%
  \tagmcend%
  \tagstructend%
}
\newenvironment{examproblems}[2]{%
  \def\ProblemHeadingTag{#2}%
  \AssignTaggingSocketPlug{block/list/label}{examproblem}%
  \begin{enumerate}%
  \setcounter{enumi}{#1}%
  \renewcommand{\labelenumi}{\textbf{\theenumi.}}%
  \setlength{\topsep}{0.35em}%
  \setlength{\partopsep}{0pt}%
  \setlength{\itemsep}{0.48em}%
  \setlength{\parsep}{0.18em}%
}{\end{enumerate}}
\newenvironment{examsubparts}{%
  \AssignTaggingSocketPlug{block/list/label}{default}%
  \begin{enumerate}%
  \setlength{\topsep}{0.18em}%
  \setlength{\partopsep}{0pt}%
  \setlength{\itemsep}{0.16em}%
  \setlength{\parsep}{0.1em}%
}{\end{enumerate}}
\newenvironment{examinstructions}{%
  \par\begingroup\itshape
}{%
  \par\endgroup\vspace{0.18em}
}
\makeatletter
\renewcommand\subsection{\@startsection{subsection}{2}{0pt}%
  {-1.25ex plus -.25ex minus -.1ex}%
  {0.55ex plus .1ex}%
  {\normalfont\large\bfseries}}
\renewcommand{\@maketitle}{%
  \newpage\null\vskip -1em%
  {\footnotesize\bfseries
    \href{https://www.ufl.edu/}{University of Florida}, \href{https://math.ufl.edu/}{Department of Mathematics}\par}%
  \vskip -0.5em%
  \tagstructbegin{tag=H1,title={Algebra First-Year Exam, January 2023, Part 2}}%
  \tagpdfparaOff%
  \tagmcbegin{tag=H1}%
    {\LARGE\bfseries\href{https://gma.math.ufl.edu/exams/algebra-first-year/}{Algebra First-Year Exam}, January 2023, Part 2\par}%
  \tagmcend%
  \tagpdfparaOn%
  \tagstructend%
  \vskip 0.25em%
  \tagmcbegin{artifact}%
    \hrule height 1pt%
  \tagmcend%
  \vskip 1em%
}
\makeatother
\title{Algebra First-Year Exam, January 2023, Part 2}
\author{}
\date{}
\begin{document}
\maketitle
\thispagestyle{empty}
\begin{examinstructions}
Answer 4 questions. Write your solutions clearly in complete English sentences. You may quote results (within reason) as long as you state them clearly.
\end{examinstructions}
\begin{examproblems}{0}{H2}
\def\ProblemHeadingText{Problem 1.}
\item
Prove that the following polynomials are irreducible in \(\mathbb{Q}[x]\), stating clearly any theorems you wish to apply.
\begin{examsubparts}
\item[\textnormal{(a)}]
\(x^3-x^2+x+5\).
\item[\textnormal{(b)}]
\(x^6-300\).
\item[\textnormal{(c)}]
\(x^4-4x^3+2x^2+21x-7\).
\end{examsubparts}
\def\ProblemHeadingText{Problem 2.}
\item
Show that the ring \(\mathbb{Z}[\sqrt{-2}]\) is a Euclidean domain with respect to the norm \(N(a+b\sqrt{-2})=a^2+2b^2\). (You may regard this ring as a subring of the field of complex numbers, with \(\sqrt{-2}=i\sqrt{2}\).)
\def\ProblemHeadingText{Problem 3.}
\item
Let~\(R\) and~\(S\) be commutative rings with 1 and \(\phi:R\to S\) a ring homomorphism with \(\phi(1)=1\). Let~\(I\) be an ideal of~\(S\) and \(\phi^{-1}(I)=\{r\in R\mid \phi(r)\in I\}\).
\begin{examsubparts}
\item[\textnormal{(a)}]
Show that \(\phi^{-1}(I)\) is an ideal of~\(R\).
\item[\textnormal{(b)}]
Show that if~\(I\) is a prime ideal then so is \(\phi^{-1}(I)\).
\item[\textnormal{(c)}]
If~\(I\) is maximal, is \(\phi^{-1}(I)\) necessarily maximal? Give a proof or counterexample.
\end{examsubparts}
\def\ProblemHeadingText{Problem 4.}
\item
Let~\(R\) be a ring with 1 and~\(M\) a (unital) left~\(R\)-module.
\begin{examsubparts}
\item[\textnormal{(a)}]
State what it means for~\(M\) to satisfy the Ascending Chain Condition (i.e. for~\(M\) to be Noetherian).
\item[\textnormal{(b)}]
Prove that if~\(M\) satisfies the ACC then~\(M\) is finitely generated.
\item[\textnormal{(c)}]
Give an example of a finitely generated~\(R\)-module~\(M\) that does not satisfy the ACC.
\end{examsubparts}
\def\ProblemHeadingText{Problem 5.}
\item
Use the rational canonical form to find one representative of each conjugacy class of elements of order 6 in \(\mathrm{GL}(4,\mathbb{Q})\).
\def\ProblemHeadingText{Problem 6.}
\item
Let \(K\subset F\subset E\) be fields.
\begin{examsubparts}
\item[\textnormal{(a)}]
If \(|F:K|=n\) and \(|E:F|=m\), where~\(m\) and~\(n\) are finite, show that \(|E:K|=nm\).
\item[\textnormal{(b)}]
Prove that if \(|E:K|\) is finite, then \(|F:K|\) and \(|E:F|\) are both finite.
\end{examsubparts}
\end{examproblems}
\end{document}
